% ============================================
% Appendix B: Protocol Finite State Machines
% ============================================
\label{app:fsm}

We provide the complete finite state machines for the three most complex
protocol interactions.

\subsection{Election State Machine}

\begin{figure}[h]
  \centering
  \begin{tikzpicture}[
    font=\footnotesize,
    >=Stealth,
    state/.style={draw, rounded corners, minimum width=1.8cm, minimum height=0.6cm, align=center, fill=gray!10},
    trans/.style={->, thick},
  ]
    \node[state] (follower) {\footnotesize FOLLOWER};
    \node[state, below left=1.2cm and 0.8cm of follower] (candidate) {\footnotesize CANDIDATE};
    \node[state, below right=1.2cm and 0.8cm of follower] (head) {\footnotesize HEAD};
    
    \draw[trans] (follower) -- (candidate) node[midway, left, font=\tiny] {head timeout};
    \draw[trans] (candidate) -- (head) node[midway, right, font=\tiny] {score wins};
    \draw[trans] (candidate) -- (follower) node[midway, above right, font=\tiny] {higher score};
    \draw[trans] (head) -- (follower) node[midway, right, font=\tiny] {higher score join};
    \draw[trans, loop above] (follower) to (follower);
    \draw[trans, loop below] (head) to (head);
  \end{tikzpicture}
  \caption{Cluster-head election FSM. All nodes start as FOLLOWER.
  On head heartbeat timeout, transition to CANDIDATE and broadcast ELECTION.
  If self.score is highest, become HEAD. If a head with higher score appears,
  demote to FOLLOWER.}
  \label{fig:election_fsm}
\end{figure}

\subsection{Submap Transfer FSM (Sender Side)}

\begin{figure}[h]
  \centering
  \begin{tikzpicture}[
    font=\footnotesize,
    >=Stealth,
    state/.style={draw, rounded corners, minimum width=1.5cm, minimum height=0.55cm, align=center, fill=gray!10},
    trans/.style={->, thick},
  ]
    \node[state] (idle) {\footnotesize IDLE};
    \node[state, right=1.5cm of idle] (tx) {\footnotesize TX\_ACTIVE};
    \node[state, right=1.5cm of tx] (wait) {\footnotesize WAIT\_ACK};
    
    \draw[trans] (idle) -- (tx);
    \draw[trans] (tx) -- (wait);
    \draw[trans] (wait) -- (idle);
    \draw[trans] (wait) edge[loop above] (wait);
    \draw[trans] (tx) edge[loop below] (tx);
  \end{tikzpicture}
  \caption{Submap transfer FSM (sender side). On SUBMAP\_NACK, re-send missing
  chunks up to 3 times; on 3rd failure, mark submap as failed and notify
  application layer to re-extract at higher compression.}
  \label{fig:submap_fsm}
\end{figure}

\subsection{Connection Lifecycle}

The connection lifecycle has five phases:

\begin{enumerate}
  \item \textbf{Discovery (0--500 ms).} Node broadcasts HELLO with signed
    NodeDescriptor. Existing nodes reply; head sends cluster digest.
  \item \textbf{Calibration handshake (0--2 s).} If $W_i$ unaligned, head sends
    coarse $T_{W_h\leftarrow W_i}$ prior. Refined after first submap exchange.
  \item \textbf{Steady state (indefinite).} POSE\_STREAM at 20 Hz,
    SUBMAP\_PUSH on boundaries, KEYFRAME at 1 Hz, MAP\_DELTA from head at 1 Hz.
  \item \textbf{Autonomous mode (on link loss).} HEARTBEAT timeout over 3 s:
    buffer submaps locally, broadcast HELLO every 1 s. On reconnect, flush
    buffered submaps in chronological order.
  \item \textbf{Leave (on exit).} Send LEAVE; head marks node's submaps as
    owned-but-absent, kept time $T_{\text{keep}}=60$ s, then garbage collected
    unless referenced by loop constraints.
\end{enumerate}

\subsection{Map Consistency Guarantees}

Co-LIC2 provides the following consistency guarantees under the protocol:

\begin{itemize}
  \item \textbf{Monotonic map version.} MAP\_DELTA carries sequential version
    numbers. A member only applies delta $v$ if its local version is $v-1$.
    Lagging members pull full snapshots.
  \item \textbf{Submap idempotency.} Globally unique submap IDs prevent
    duplicate merge. The head maintains a set of recently merged IDs and
    discards duplicates.
  \item \textbf{PGO convergence.} iSAM2 incrementally updates the graph;
    PGO corrections are applied as rigid submap transforms, guaranteeing
    that a submap's internal geometry is preserved.
  \item \textbf{Orphan timeout.} Absent-node submaps are kept for 60 s,
    then garbage collected. If a loop constraint references an orphan
    submap, it is preserved until the constraint is removed.
\end{itemize}
